\subsection{Coxeter 变换}

设 C 是 R 的一个室，\(\{\alpha_{1},\ldots,\alpha_{l}\}\) 是相应的 R 的基，并令 \(c=s_{\alpha_{1}}\ldots s_{\alpha_{l}}\)。W 的元素 c 称为由 C 和双射 \(i\mapsto\alpha_{i}\) 定义的 W 的 Coxeter 变换（Chap. V, § 6, no. 1）。它的阶 h 称为 W（或 R）的 Coxeter 数。

命题 30. 假设 R 不可约。设 m 是介于 1 和 h - 1 之间且与 h 互素的整数。则 \(\exp\left(\frac{2i\pi m}{h}\right)\) 是 c 的一个重数为 1 的特征值。

特别地，\(m\) 是 \(W\) 的一个指数，cf.~Chap. V, § 6, no. 2。

我们先证明一个引理：

引理 4. 对所有 \(w \in W\)，\(w\) 的特征多项式具有整数系数。

我们知道 \(\{\alpha_{1},\ldots,\alpha_{l}\}\) 是由 R 生成的 V 的子群 \(Q(R)\) 的一个基（no. 6, Th. 3）。由于 w 保持 \(Q(R)\) 不变，w 关于 \(\{\alpha_{1},\ldots,\alpha_{l}\}\) 的矩阵具有整数项；因此其特征多项式具有整数系数。

令 P 为 c 的特征多项式。上述引理表明 P 的系数是整数。根据 Chap. V, § 6, no. 2, Cor. 2 of Prop. 3，h 次本原单位根 \(z = \exp\left(\frac{2i\pi}{h}\right)\) 是 P 的一个单根。因此，z 在 Q 上的每个共轭元也是 P 的一个单根。但我们知道（Algebra, Chap. V），h 次本原单位根在 Q 上共轭。因此它们全都是 P 的单根，这就证明了该命题。

命题 31. 假设 R 不可约且约化，且 \(\beta = n_{1}\alpha_{1} + \cdots + n_{l}\alpha_{l}\) 是 R 的最高根（cf.~no. 8）。则 \(n_{1} + \cdots + n_{l} = h - 1\)。

令 \(R_{+}\) 为关于 C 的正根集合。于是（no. 10, Cor. of Prop. 29）：

\[
\begin{array}{l} n _ {1} + \dots + n _ {l} = \frac {1}{2} \sum_ {\alpha \in \mathrm{R} _ {+}} n (\beta , \alpha) \\ = 1 + \frac {1}{2} \sum_ {\alpha \in \mathrm{R} _ {+}, \alpha \neq \beta} n (\beta , \alpha) = 1 + \sum_ {\alpha \in \mathrm{R} _ {+}, \alpha \neq \beta} \frac {(\alpha | \beta)}{(\alpha | \alpha)}. \end{array}
\]

由第 8 号的命题 25 (iv)，对所有 \(\alpha \in \mathbb{R}_+\) 且 \(\alpha \neq \beta\)，\(n(\alpha, \beta) = 0\) 或 1，所以 \(n(\alpha, \beta)^2 = n(\alpha, \beta)\)，即 \(\frac{4(\alpha | \beta)^2}{(\beta | \beta)^2} = \frac{2(\alpha | \beta)}{(\beta | \beta)}\)。于是：

\[
\begin{array}{l} n _ {1} + \dots + n _ {l} + 1 = 2 + 2 \sum_ {\alpha \in R _ {+}, \alpha \neq \beta} \frac {(\alpha | \beta) ^ {2}}{(\alpha | \alpha) (\beta | \beta)} \\ = 2 \sum_ {\alpha \in R _ {+}} \frac {(\alpha | \beta) ^ {2}}{(\alpha | \alpha) (\beta | \beta)} = (\beta | \beta) ^ {- 1} \sum_ {\alpha \in R} (\frac {\alpha}{\| \alpha \|} | \beta) ^ {2}. \end{array}
\]

由 Chap. V, § 6, no. 2 的 Cor. of Th. 1，

\[
\sum_ {\alpha \in R} \left(\frac {\alpha}{\| \alpha \|} | \beta\right) ^ {2} = h (\beta | \beta)
\]

\[
\mathrm{so} n _ {1} + \dots + n _ {l} + 1 = h.
\]

命题 32. 假设 R 不可约，且所有根具有相同长度。设 \(\alpha \in R\)。与 \(\alpha\) 不正交的 R 的元素数目为 4h - 6。

令 \(R'\) 为不与 \(\alpha\) 成比例且不与其正交的根的集合。根据 Chap. V, § 6, no. 2, Cor. of Th. 1，

\[
(\alpha | \alpha) ^ {2} + (\alpha | - \alpha) ^ {2} + \sum_ {\beta \in \mathrm{R} ^ {\prime}} (\alpha | \beta) ^ {2} = h (\alpha | \alpha) ^ {2},
\]

即

\[
\sum_ {\beta \in \mathsf {R} ^ {\prime}} (\alpha \mid \beta) ^ {2} = (h - 2) (\alpha \mid \alpha) ^ {2}.
\]

若 \(\beta \in \mathbf{R}'\)，则由第 3 号的列表，\((\alpha | \beta) = \pm \frac{1}{2} (\alpha | \alpha)\)。于是

\[
\frac {1}{4} \operatorname{Card} R ^ {\prime} = h - 2, \quad \operatorname{Card} R ^ {\prime} = 4 h - 8,
\]

与 \(\alpha\) 不正交的根的数目为 \(\operatorname{Card} R' + 2 = 4h - 6\)。

命题 33. 假设 R 不可约且约化。置 \(s_{\alpha_i} = s_i\)，并令 \(\Gamma\) 为由 \(c = s_1 \ldots s_l\) 生成的 W 的子群。

\begin{enumerate}
\def\labelenumi{(\roman{enumi})}
\item
  令 \(\theta_{i} = s_{l}s_{l - 1}\dots s_{i + 1}(\alpha_{i})\)（\(i = 1,\dots ,l\)）。于是 \(\theta_{i} > 0,c(\theta_{i}) < 0\)。
\item
  若 \(\alpha\) 是满足 \(>0\) 且 \(c(\alpha) < 0\) 的根，则 \(\alpha\) 等于某个 \(\theta_{i}\)。
\item
  族 \((\theta_{i})_{1\leqslant i\leqslant l}\) 是 \(\mathrm{Q}(\mathrm{R})\) 的一个基。
\item
  令 \(\Omega_{i}\) 为 \(\theta_{i}\) 在 \(\Gamma\) 下的轨道。各集合 \(\Omega_{i}\) 两两不交，它们是 \(\Gamma\) 在 \(\mathbb{R}\) 上的全部轨道，并且每个集合含有 \(h\) 个元素。
\end{enumerate}

首先注意，\((s_1, \ldots, s_l)\) 是 \(c\) 的一个约化分解（Chap. IV, § 1, no. 1），相对于集合 \(S\)（由 \(s_i\) 构成）。事实上，否则会存在一个子集 \(X = S - \{j\}\)，它含有 \(l - 1\) 个 \(S\) 的元素，使得 \(c \in W_X\)，这将与 Chap. IV, § 1, no. 8 的 Prop. 7 的 Cor. 2 矛盾。

将第 6 号的 Prop. 17 的 Cor. 2 应用于 \(c\)，得到断言 (i) 和 (ii)。

令 \(Q_{i}\) 为 \(Q(R)\) 的子群，由 \(\alpha_{j}, j > i\) 生成。显然 \(Q_{i}\) 在 \(s_{j}, j > i\) 下保持不变，并且 \(s_{j}(\alpha_{i}) = \alpha_{i} \mod Q_{i}\) 对 \(j > i\) 成立。因此：

\[
\theta_ {i} = s _ {l} \dots s _ {i + 1} (\alpha_ {i}) = \alpha_ {i} \bmod . Q _ {i}.
\]

换言之，存在整数 \(c_{ij}\)，使得

\[
\theta_ {i} = \alpha_ {i} + \sum_ {j > i} c _ {i j} \alpha_ {j}.
\]

(iii) 立即得到。

最后，设 \(\alpha\) 是一个根。元素 \(\sum_{k=0}^{h-1} c^k(\alpha)\) 在 \(c\) 下不变，因而为零（Chap. V, § 6, no. 2）。因此，\(c^k(\alpha)\) 不可能全都同号，存在一个 \(k\)，使得 \(c^k(\alpha) > 0\) 且 \(c^{k+1}(\alpha) < 0\)。由 (ii)，\(c^k(\alpha)\) 是某个 \(\theta_i\)。因此，\(\Gamma\) 在 R 上的每个轨道都是某个 \(\Omega_i\)。将 \((x|y)\) 延拓为 V \(\otimes\) C 上的厄米型。根据上述注记，\(\Gamma\) 在 R 中的每个轨道至多有 \(h\) 个元素，且至多有 \(l\) 个不同的轨道。现在（Chap. V, § 6, no. 2, Theorem 2, ii)），R 的基数等于 \(hl\)，这立即推出 (iv)。
% semantic-fragments: legacy-input-seam
\relax{}
